% ===================================================================
%  TABELUL Nr. 10 — Marea. — Portul.
%  Traduction 2026 d'apres texto/io, controlee sur texto/fr, texto/en
%  et texto/es. Consigne : voir texto/ro/10-tabelul-01.tex.
%
%  Le francais decale sa numerotation d'un cran a partir de l'alinea 8.
%  C'est une coquille de l'imprimeur francais ; le roumain suit l'ido.
%
%  LE TABLEAU DE LA MER EST LE PLUS EMPRUNTE DE LA COLONNE, comme en
%  ukrainien et en basque, et pour la meme raison : la marine moderne
%  est arrivee toute faite au XIXe siecle. « cuirasat », « torpilor »,
%  « fregata », « doc », « elice », « pilot » sont les memes mots
%  partout. Ce qui reste roumain, ce sont les mots de la terre :
%  « țărm » le rivage, « mal » la rive, « stânca » le rocher.
% ===================================================================

%%K t10-apar-1 apar
\VUsaut{4.0mm}
\VUtitre{15.0pt}{60}{TABELUL Nr. 10}
\VUsaut{3.0mm}
\VUpk{11.4pt}{40}{Marea. --- Portul.}
\VUsaut{3.0mm}

%%K t10-01-1 p
1. --- Vara, în vreme ce unii caută liniștea și răcoarea la munte, alții preferă să
meargă la mare. Așa am făcut și noi anul trecut, fiindcă iubim mult
\VUgras{oceanul} \textsuperscript{(1)}.

%%K t10-02-1 p
2. --- Cât despre mine, îmi face mare plăcere să privesc această întindere de apă
fără margini, așternută până la \VUgras{zare} \textsuperscript{(2)}, și să văd cum
uriașele \VUgras{vapoare} \textsuperscript{(3)} pleacă în \VUgras{cursă} lungă cu
călătorii care merg în America.

%%K t10-03-1 p
3. --- De aceea tata, care îmi știe gusturile, alege totdeauna pentru vacanță un
țărm foarte liniștit lângă unul dintre marile noastre \VUgras{porturi militare}.

%%K t10-03-2 p
La ultima noastră venire \VUgras{escadra} a sosit și a ancorat în spatele uriașului
\VUgras{dig} \textsuperscript{(4)}, care închide și apără \VUgras{rada militară}
\textsuperscript{(5)}, și am putut să admir pe îndelete feluritele \VUgras{nave de
război}: de la micul \VUgras{submarin} \textsuperscript{(10)}, care se cufundă și
piere sub apă, până la uriașul \VUgras{cuirasat} \textsuperscript{(9)}, care până
acum nu avea de ce se teme decât de atacurile \VUgras{torpilorului}
\textsuperscript{(8)}.

%%K t10-tit-2 sub
\VUsaut{3.0mm}
\VUpk{10.2pt}{40}{Navele mici.}
\VUsaut{3.0mm}

%%K t10-04-1 p
4. --- Acesta din urmă știa deja să găsească \VUgras{cu îndemânare} locul slab din
platoșa groasă de metal și să pună acolo \VUgras{torpila} primejdioasă, a cărei
explozie scufundă nava; dar totuși trebuia să fie temut mai puțin decât submarinul,
fiindcă distrugătoarele se iau ușor după el.

%%K t10-05-1 p
5. --- Am putut vedea și \VUgras{crucișătoarele} cuirasate iuți
\textsuperscript{(6)}, aproape tot atât de nevătămabile ca un adevărat cuirasat,
câteva \VUgras{transportoare} \textsuperscript{(12)}, trei sau patru
\VUgras{canoniere} \textsuperscript{(7)} și, în sfârșit, \VUgras{fregata}
\textsuperscript{(11)}. \VUgras{Cea mai frumoasă priveliște este a acestei flote
care ridică ancora.}

%%K t10-06-1 p
6. --- Ce deosebire de alcătuire între acești munți de oțel și elegantele
\VUgras{corăbii cu trei catarge} sau grațioasele \VUgras{veliere}
\textsuperscript{(13)}, care se mai folosesc în flota comercială și pe care le-am
văzut intrând în port cu toate pânzele întinse, în vreme ce mă plimbam pe
\VUgras{dig} \textsuperscript{(14)}. E drept că pierdeau mult când \VUgras{vântul}
le era împotrivă. Trebuiau să strângă toate pânzele ca să intre în doc, și era
nevoie de un \VUgras{remorcher} \textsuperscript{(16)} ca să le ia și să le aducă
la chei, ca pe niște simple \VUgras{șlepuri} \textsuperscript{(17)}.

%%K t10-tit-3 sub
\VUsaut{3.0mm}
\VUpk{10.2pt}{40}{Portul comercial.}
\VUsaut{3.0mm}

%%K t10-07-1 p
7. --- Portul despre care vorbesc este interesant tocmai prin aceea că are nu numai
rada militară, făcută în chip meșteșugit prin lucrări uriașe, ci și un minunat
\VUgras{port comercial} la \VUgras{gura} unui râu mare.

%%K t10-08-1 p
8. --- Fâșia de pământ care desparte cele două docuri este întărită și apărată de
un șir de \VUgras{baterii} și de \VUgras{bastioane} împinse în mare. Ca să te
plimbi pe \VUgras{metereze} \textsuperscript{(15)}, trebuie o învoire deosebită. Am
căpătat-o ușor prin unchiul meu, care este inginer la \VUgras{șantierele navale}
\textsuperscript{(18)}, și am mers să văd cum se construiesc navele sub șoproane
uriașe.

%%K t10-09-1 p
9. --- Am văzut chiar lansarea unui cuirasat și îmi amintesc tulburarea care a
cuprins toată mulțimea când uriașul munte, lăsat în voia lui, a lunecat de pe
patul său și a ieșit la fața apei râului.

%%K t10-10-1 p
10. --- Ca să ajungi la șantiere, treci \VUgras{piața de arme}
\textsuperscript{(19)}, unde se instruiesc în fiecare zi trupele garnizoanei, și
mergi pe \VUgras{podețul} de fier \textsuperscript{(20)}, care leagă piața de
\VUgras{arsenal} \textsuperscript{(21)}.

%%K t10-tit-4 sub
\VUsaut{3.0mm}
\VUpk{10.2pt}{40}{Arsenalul și docurile.}
\VUsaut{3.0mm}

%%K t10-11-1 p
11. --- Cu unchiul am cercetat acest așezământ mare, plin de \VUgras{tunuri}
\textsuperscript{(22)}, de \VUgras{ghiulele} \textsuperscript{(23)} și de
\VUgras{proiectile}. Am văzut și \VUgras{cazangeria} \textsuperscript{(24)}, unde
se fac \VUgras{cazanele} \textsuperscript{(25)} pentru vapoare.

%%K t10-12-1 p
12. --- Chiar alături sunt \VUgras{docurile} \textsuperscript{(26)}, docuri uscate,
și foarte deosebitele \VUgras{docuri plutitoare} \textsuperscript{(27)}, în care am
putut cerceta de aproape, pe o navă în reparație, \VUgras{elicea}
\textsuperscript{(28)}, \VUgras{chila} \textsuperscript{(29)} și \VUgras{cârma}
\textsuperscript{(30)} corăbiei.

%%K t10-13-1 p
13. --- Portul comercial merită să fie cercetat nu mai puțin decât cel militar, și
am petrecut multe ceasuri hoinărind pe chei, unde stau \VUgras{vapoarele cu zbaturi}
\textsuperscript{(31)}, care urcă pe râu, și privind cum lucrează \VUgras{macaraua
rotitoare} \textsuperscript{(32)}, ridicând greutăți uriașe ca pe niște fulgi.

%%K t10-tit-5 sub
\VUsaut{4.0mm}
\VUpk{10.2pt}{40}{Pilotul.}
\VUsaut{3.0mm}

%%K t10-14-1 p
14. --- Admiram \VUgras{transatlanticele} \textsuperscript{(34)}, care părăseau
rada cu \VUgras{steagurile} fâlfâind \textsuperscript{(35)}, conduse de un
\VUgras{pilot} iscusit \textsuperscript{(36)}, care știa minunat să se țină pe
\VUgras{șenal}, însemnat cu \VUgras{geamanduri} \textsuperscript{(40)} și cu
\VUgras{balize}, ocolind \VUgras{barjele} \textsuperscript{(41)}, atât de greu de
cârmuit sub singura lor pânză pătrată. Deosebeam la \VUgras{proră}
\textsuperscript{(37)} \VUgras{cabinele} \textsuperscript{(39)} de clasa a doua,
iar la \VUgras{pupă} \textsuperscript{(38)} --- saloanele pentru călătorii de clasa
întâi.

%%K t10-15-1 p
15. --- Uneori priveam și încărcarea unui \VUgras{cargobot}
\textsuperscript{(42)}, în care se vărsau mărfuri din \VUgras{depozit}
\textsuperscript{(43)} și din \VUgras{gara portului} \textsuperscript{(44)}, aduse
de trenurile numeroase ale \VUgras{liniei de chei} \textsuperscript{(45)}.

%%K t10-tit-6 sub
\VUsaut{3.0mm}
\VUpk{10.2pt}{40}{Plimbarea cu barca.}
\VUsaut{3.0mm}

%%K t10-16-1 p
16. --- Mă plimbam adesea cu barca în \VUgras{docul închis}
\textsuperscript{(53)}, unde se adăpostesc \VUgras{bărcile} mici
\textsuperscript{(46)}, și era o bucurie pentru mine să țin o \VUgras{vâslă}
\textsuperscript{(47)}. Urcam pe \VUgras{puntea} \textsuperscript{(48)}
\VUgras{bărcii cu pânze} \textsuperscript{(49)}, mă cățăram pe
\VUgras{parâme} \textsuperscript{(52)} până la \VUgras{catarg}
\textsuperscript{(51)}, până la \VUgras{vergi} \textsuperscript{(50)}, iar apoi
îmi urmam plimbarea într-o \VUgras{șalupă} \textsuperscript{(54)} printre grelele
\VUgras{vase de pescuit} \textsuperscript{(55)}, unde se usucă \VUgras{plasele}
întinse \textsuperscript{(56)}.

%%K t10-17-1 p
17. --- Pe \VUgras{chei} \textsuperscript{(58)} treceau prin fața mea cele mai
felurite \VUgras{mărfuri} \textsuperscript{(59)}, împachetate în \VUgras{lăzi}
\textsuperscript{(60)} sau în \VUgras{baloturi} \textsuperscript{(61)}. Alcătuiau
\VUgras{încărcătura} \textsuperscript{(63)} șlepurilor, iar \VUgras{macarale}
puternice \textsuperscript{(62)} le scoteau și le puneau pe \VUgras{camioane}
\textsuperscript{(64)}, în vreme ce \VUgras{cărăușii} le declarau și plăteau vama la
\VUgras{vamă} \textsuperscript{(65)}.

%%K t10-18-1 p
18. --- În sfârșit, ajungând la \VUgras{podul mobil} \textsuperscript{(66)} sau la
\VUgras{ecluză} \textsuperscript{(69)}, trecând pe lângă \VUgras{draga}
\textsuperscript{(68)}, care curăță \VUgras{canalul} \textsuperscript{(67)},
acostam la dig lângă \VUgras{catargul de semnale} \textsuperscript{(70)}.

%%K t10-19-1 p
19. --- Tocmai de cealaltă parte a acestui dig acostează \VUgras{șalupele}
\textsuperscript{(71)}, care aduc la țărm ofițerii escadrei. Este frumos să vezi
cum ridică marinarii vâslele deodată, iar barca, purtată de \VUgras{val}
\textsuperscript{(72)}, se apropie ușor de scară. De aici se văd cel mai bine
\VUgras{fluxul} și mișcarea \VUgras{refluxului} și a \VUgras{resacului} pe
\VUgras{plajă} \textsuperscript{(73)}.

%%K t10-tit-7 sub
\VUsaut{3.0mm}
\VUpk{10.2pt}{40}{Cercul ofițerilor.}
\VUsaut{3.0mm}

%%K t10-20-1 p
20. --- Ca să mă întorc acasă, mă suiam în \VUgras{tramvaiul electric}
\textsuperscript{(74)}, a cărui \VUgras{stație} \textsuperscript{(75)} este lângă
\VUgras{stâlpul} din fața cercului ofițerilor.

%%K t10-20-2 p
De pe \VUgras{balconul} \textsuperscript{(76)} cercului, împodobit cu
\VUgras{ancore} \textsuperscript{(77)}, cu tunuri și cu ghiulele, vedeam adesea
societatea strălucită a ofițerilor de marină: \VUgras{locotenenți}
\textsuperscript{(78)} cu spade, \VUgras{căpitani de artilerie}
\textsuperscript{(79)} și de \VUgras{infanterie} \textsuperscript{(80)} cu
\VUgras{săbii}. În spatele lor stătea un \VUgras{marinar}
\textsuperscript{(81)}, care le dădea \VUgras{luneta} când voiau să cerceteze ce se
petrece în depărtare.

%%K t10-21-1 p
21. --- Am văzut și tineri \VUgras{aspiranți} \textsuperscript{(82)}, care primeau
porunci de la \VUgras{comandantul} lor \textsuperscript{(83)} sau de la
\VUgras{amiral} \textsuperscript{(84)}, în vreme ce \VUgras{cartierul-maistru}
\textsuperscript{(85)} aștepta să îi ducă pe navă.

%%K t10-22-1 p
22. --- De pe acest balcon priveliștea este minunată: stăpânește toată plaja cu
\VUgras{corturile} ei \textsuperscript{(86)} și cu \VUgras{cabinele de baie}
\textsuperscript{(87)}, iar la ceasul scăldatului se vede cum \VUgras{scăldătorii}
\textsuperscript{(88)} zburdă veseli înaintea \VUgras{cazinoului}
\textsuperscript{(89)}.

%%K t10-23-1 p
23. --- La capătul plajei țărmul face o mică \VUgras{peninsulă} stâncoasă
\textsuperscript{(91)}, de care se sfarmă \VUgras{talazurile}
\textsuperscript{(92)} și pe care este zidit \VUgras{farul} \textsuperscript{(93)},
care noaptea arată marinarilor drumul. Această peninsulă este legată de uscat numai
printr-un \VUgras{istm} foarte îngust \textsuperscript{(90)}.

%%K t10-tit-8 sub
\VUsaut{3.0mm}
\VUpk{10.2pt}{40}{Priveliștea de obște a țărmului.}
\VUsaut{3.0mm}

%%K t10-24-1 p
24. --- \VUgras{Faleza} \textsuperscript{(94)} se întinde mai departe, tăiată de
mare, care sapă în ea un șir ciudat de \VUgras{golfuri} \textsuperscript{(95)} și
de \VUgras{capuri}, făcând-o foarte pitorească.

%%K t10-24-2 p
Pe \VUgras{podișul} \textsuperscript{(96)} care stăpânește \VUgras{strâmtoarea}
\textsuperscript{(98)} stau un \VUgras{fort} foarte însemnat \textsuperscript{(97)}
și câteva vile.

%%K t10-25-1 p
25. --- Această strâmtoare desparte de uscat o \VUgras{insulă} foarte primejdioasă
\textsuperscript{(99)}, înconjurată de \VUgras{stânci} \textsuperscript{(100)} și
de \VUgras{recifuri}, unde pier uneori bărci și chiar nave mari. Oricât de sprintenă
ar fi ajutorarea și oricât de repede ar veni \VUgras{bărcile de salvare}, aceste
\VUgras{naufragii} sunt înspăimântătoare, iar în \VUgras{furtunile} de echinocțiu
câteva nave au pierit cu tot echipajul și cu toată încărcătura.

%%K t10-25-2 p
Cu tot acest vecin, \VUgras{portul} acesta este cercetat bucuros de marinari.
\VUsaut{6.0mm}
\VUfilet{22mm}
